\section{Conclusion and Future Work}\label{sec:Conclusion}

We introduced and studied the problem of finding degree- and h-index-based hubs on unmaterialized relational graphs (\ced). We proposed a sketch propagation framework for approximate degree-based \ced queries and extended it to h-index-based \ced queries with a pivot-based algorithm. We also considered personalized \ced queries for both centrality measures and enhanced efficiency with early termination. Theoretically, our sketch propagation and pivot-based algorithms answer approximate \ced queries correctly with high probability. We verified the effectiveness and efficiency of our proposed algorithms with extensive experimentation on real-world heterogeneous data.

For future work, it would be interesting to extend our framework to process hub queries over relational graphs induced by general meta-graphs, which impose substantial technical challenges due to the potentially complex structure of meta-graphs. Another possible direction involves generalizing sketch-based methods for graph analysis tasks over relational graphs other than hub queries~\cite{niu2025sans}.

\section*{Acknowledgment}
This research is partly supported by the Lee Kong Chian Fellowship of Singapore Management University.
P. Karras was supported by the Independent Research Fund Denmark (Project 9041-00382B).
Y. Wang was supported by the National Natural Science Foundation of China (Grant No.~62202169).